\documentclass[12pt,a4paper]{article}

% -------------------------------------------------
% Encoding and language
% -------------------------------------------------
\usepackage[utf8]{inputenc}
\usepackage[T1]{fontenc}
\usepackage[english]{babel}

% -------------------------------------------------
% Page layout
% -------------------------------------------------
\usepackage[
a4paper,
top=2.2cm,
bottom=2.2cm,
left=2.5cm,
right=2.5cm
]{geometry}

% -------------------------------------------------
% Typography
% -------------------------------------------------
\usepackage{lmodern}
\usepackage{microtype}

% -------------------------------------------------
% Mathematics
% -------------------------------------------------
\usepackage{amsmath}
\usepackage{amssymb}

% -------------------------------------------------
% Tables and lists
% -------------------------------------------------
\usepackage{array}
\usepackage{longtable}
\usepackage{booktabs}
\usepackage{enumitem}

% -------------------------------------------------
% Colors
% -------------------------------------------------
\usepackage{xcolor}

\definecolor{primary}{RGB}{31,78,121}
\definecolor{secondary}{RGB}{68,68,68}
\definecolor{lightgray}{RGB}{245,245,245}

% -------------------------------------------------
% Headers and footers
% -------------------------------------------------
\usepackage{fancyhdr}

\pagestyle{fancy}
\fancyhf{}

\fancyhead[L]{\small Periodic Teaching Plan}
\fancyhead[R]{\small PGF-02-R39}

\fancyfoot[C]{\thepage}

\renewcommand{\headrulewidth}{0.4pt}
\renewcommand{\footrulewidth}{0pt}

% -------------------------------------------------
% Section formatting
% -------------------------------------------------
\usepackage{titlesec}

\titleformat{\section}
{\Large\bfseries\color{primary}}
{\thesection.}
{0.5em}
{}

\titleformat{\subsection}
{\large\bfseries\color{primary}}
{}
{0pt}
{}

\titleformat{\subsubsection}
{\normalsize\bfseries\color{secondary}}
{}
{0pt}
{}

% -------------------------------------------------
% Links
% -------------------------------------------------
\usepackage[
colorlinks=true,
linkcolor=primary,
urlcolor=primary,
citecolor=primary
]{hyperref}

% -------------------------------------------------
% Paragraph formatting
% -------------------------------------------------
\setlength{\parindent}{0pt}
\setlength{\parskip}{0.6em}

% -------------------------------------------------
% List formatting
% -------------------------------------------------
\setlist[itemize]{
	leftmargin=1.5cm,
	itemsep=0.2em,
	topsep=0.2em
}

\setlist[description]{
	leftmargin=1.5cm,
	labelindent=0cm,
	style=nextline,
	itemsep=0.5em
}

% -------------------------------------------------
% Custom commands
% -------------------------------------------------

\newcommand{\criterio}[2]{
	\item[\textbf{#1:}] #2
}

\newcommand{\placeholder}[1]{
	\textit{\textcolor{secondary}{[#1]}}
}

% -------------------------------------------------
% Document
% -------------------------------------------------

\begin{document}
	
	% =================================================
	% TITLE
	% =================================================
	
	\begin{center}
		
		{\Large\bfseries\color{primary}
			PERIODIC TEACHING PLAN}
		
		\vspace{0.4cm}
		
		{\normalsize\bfseries PGF-02-R39}
		
		\vspace{0.3cm}
		
		{\normalsize
			\textbf{Date:} 12-08-2026 to 05-11-2026}
		
	\end{center}
	
	\vspace{0.8cm}
	
	\tableofcontents
	
	
	% =================================================
	% 1. IDENTIFICATION
	% =================================================
	
	\section{Identification}
	
	\textbf{Name of the area / skill:}
	
	Seminars -- ``Entrepreneurship and Innovation''
	
	\textbf{Seminar research question:}
	
	How can artificial intelligence help us design and build technological solutions to problems in our environment?
	
	\textbf{Teacher:}
	
	Nicolás López
	
	\textbf{Grade:}
	
	10th
	
	\textbf{Term:}
	
	37
	
	% =================================================
	% 2. GRADE-LEVEL OBJECTIVE
	% =================================================
	
	\section{Grade-Level Objective}
	
	\placeholder{Information not specified in the provided document.}
	
	% =================================================
	% 3. LEARNING OBJECTIVE
	% =================================================
	
	\section{Learning Objective}
	
	\subsection*{Learning Objective}
	
	Identify needs, problems, or opportunities for innovation in the immediate environment through contextual observation, empathy, and challenge definition.
	
	\subsection*{Assessment Criteria}
	
	\begin{description}
		
		\criterio{C1}{
			Recognizes needs, discomforts, or opportunities for improvement
			within the school or nearby community.
		}
		
		\criterio{C2}{
			Observes behaviors, spaces, users, or situations to better
			understand a need.
		}
		
		\criterio{C3}{
			Listens to potential users through conversations, guided
			questions, or empathy exercises.
		}
		
		\criterio{C4}{
			Differentiates a real need from an anticipated solution idea.
		}
		
		\criterio{C5}{
			Organizes contextual information using maps, user profiles,
			lists of findings, or simple diagrams.
		}
		
		\criterio{C6}{
			Prioritizes an innovation opportunity considering relevance,
			feasibility of action, and expected benefit.
		}
		
		\criterio{C7}{
			Defines a clear, concrete, and understandable innovation
			challenge for the group.
		}
		
		\criterio{C8}{
			Relates the challenge to knowledge from design, communication,
			economics, technology, citizenship, or another relevant
			discipline.
		}
		
		\criterio{C9}{
			Participates in initial ideation activities with openness to
			diverse proposals.
		}
		
		\criterio{C10}{
			Recognizes the importance of understanding the user before
			proposing a solution.
		}
		
	\end{description}
	
	% =================================================
	% 4. SCENARIO AND NARRATIVE
	% =================================================
	
	\section{Scenario and Narrative}
	
	\textbf{Scenario:}
	
	The context or general framework in which the learning experiences
	proposed for the term are situated. It includes the place, time,
	circumstances, problems, and other relevant elements.
	
	\textbf{Narrative:}
	
	A description or dialogic construction developed by the grade-level
	team to align the pedagogical actions.
	
	\placeholder{Specific information regarding the scenario and narrative
		was not provided in the document.}
	
	% =================================================
	% 5. CONCEPTUAL REFERENCES
	% =================================================
	
	\section{Conceptual References}
	
	\textbf{Innovation:}
	
	According to the Organisation for Economic Co-operation and Development
	(OECD) and Eurostat (2018), innovation is the process through which
	products, services, processes, or organizational methods are created
	or improved to respond more efficiently, creatively, and sustainably
	to a need or problem, generating value for people and society.
	
	Innovation generally:
	
	\begin{itemize}
		
		\item Responds to a specific need or problem.
		
		\item Provides a significant improvement over existing solutions.
		
		\item Generates value for people or society.
		
		\item May be technological, social, educational, environmental,
		or organizational.
		
		\item Requires creativity, knowledge, and the ability to implement
		ideas.
		
	\end{itemize}
	
	% =================================================
	% 6. PRE-LESSON AND CONTEXT
	% =================================================
	
	\section{Pre-Lesson and Context}
	
	A meeting will be held with the entire grade level at the Social
	Campus to explain the scope and projected development of the new
	seminar format to the students.
	
	During this meeting:
	
	\begin{itemize}
		
		\item The teaching team in charge will be introduced.
		
		\item The research questions of the different seminars will be
		presented.
		
		\item Students will be invited to check their email in the
		afternoon, upon arriving home, to complete the registration
		process.
		
		\item The logistical aspects related to the development of the
		seminars will be explained.
		
	\end{itemize}
	
	Subsequently, a special guest from CESA University will speak to the
	students about entrepreneurship, how to undertake entrepreneurial
	initiatives, and its importance for adult life.
	
	Students will be encouraged to participate through questions, opinions,
	and reflections related to entrepreneurship and innovation.
	
	% =================================================
	% 7. LEARNING EXPERIENCE
	% =================================================
	
	\section{Learning Experience}
	
	The learning experiences correspond to the pedagogical actions that
	enable the development of the learning objective while addressing the
	assessment criteria. Links to virtual resources or the platforms used
	to share the experiences (OPEN, Teams, etc.) should be included.
	
	% -------------------------------------------------
	% Pedagogical Action 1
	% -------------------------------------------------
	
	\subsection{Pedagogical Action No.\textsuperscript{o} 1}
	
	\textbf{Assessment Criteria:} C1 to C4
	
	\subsubsection*{Innovation Lab: Identifying Opportunities in My Community}
	
	Working in teams, students will conduct an observation and research
	process focused on a problem affecting the school.
	
	Possible topics include:
	
	\begin{itemize}
		
		\item Waste management.
		
		\item Mobility.
		
		\item Coexistence.
		
		\item Mental health.
		
		\item Responsible use of technology.
		
		\item Local entrepreneurship.
		
		\item Safety.
		
		\item Access to recreational spaces.
		
		\item Other problems identified by the students.
		
	\end{itemize}
	
	Students will use different information-gathering strategies, such as:
	
	\begin{itemize}
		
		\item Surveys.
		
		\item Interviews.
		
		\item Observations.
		
		\item Information search and analysis.
		
	\end{itemize}
	
	They will then analyze the causes, consequences, and impact of the
	problem and present an evidence-based diagnosis.
	
	Each team will have four classes to complete this pedagogical action.
	The teacher will continuously monitor the students' progress and
	development throughout the process.
	
	% -------------------------------------------------
	% Pedagogical Action 2
	% -------------------------------------------------
	
	\subsection{Pedagogical Action No.\textsuperscript{o} 2}
	
	\textbf{Assessment Criteria:} C5 to C10
	
	Based on the diagnosis developed previously, the teams will develop
	an innovative solution proposal that is feasible, sustainable, and
	capable of generating a positive impact on the community.
	
	The proposal may take the form of:
	
	\begin{itemize}
		
		\item A product.
		
		\item A service.
		
		\item A campaign.
		
		\item An application.
		
		\item A social strategy.
		
		\item An environmental initiative.
		
		\item A community project.
		
	\end{itemize}
	
	Throughout the process, students will be expected to:
	
	\begin{itemize}
		
		\item Justify their idea.
		
		\item Identify the resources required.
		
		\item Identify the potential beneficiaries.
		
		\item Define how they will evaluate the impact of their solution.
		
	\end{itemize}
	
	Each subject area may contribute discipline-specific elements to
	strengthen the proposals, such as:
	
	\begin{itemize}
		
		\item Business model design.
		
		\item Budget development.
		
		\item Scientific analysis.
		
		\item Effective communication.
		
		\item Visual identity.
		
		\item Prototyping.
		
		\item Sustainability.
		
		\item Citizenship skills.
		
		\item Project presentation and communication.
		
	\end{itemize}
	
	% =================================================
	% 7. LEARNING EXPERIENCE
	% =================================================
	
	\section{Learning Experience}
	
	The learning experience is organized into seven classes and two learning evidences. Classes 1--4 focus on identifying, observing, and understanding a need or innovation opportunity and culminate in Learning Evidence 1, which assesses Criteria C1--C5. Classes 5--7 focus on prioritizing the opportunity, defining an innovation challenge, and conducting initial ideation, culminating in Learning Evidence 2, which assesses Criteria C6--C10.
	
	% -------------------------------------------------
	% 7.1 CLASS 1
	% -------------------------------------------------
	
	\subsection{Class 1: Identifying Innovation Opportunities}
	
	\textbf{Assessment Criteria:} C1
	
	\textbf{Learning Purpose:}
	
	Recognize needs, discomforts, problems, and opportunities for improvement within the school or nearby community.
	
	\textbf{Main Activity:}
	
	The class begins with an introduction to the concept of innovation. Students analyze examples of innovation and distinguish between a need, a problem, a discomfort, an opportunity, and a proposed solution.
	
	Students then explore the school environment and identify situations that could represent opportunities for improvement.
	
	Working in teams, students generate an initial list of possible opportunities without immediately proposing technological or other solutions.
	
	\textbf{Student Product:}
	
	Initial opportunity map containing several possible needs, problems, or opportunities identified in the school or nearby community.
	
	\textbf{Formative Assessment:}
	
	The teacher reviews whether students are identifying actual needs or problems rather than prematurely proposing products, applications, or technologies.
	
	% -------------------------------------------------
	% 7.2 CLASS 2
	% -------------------------------------------------
	
	\subsection{Class 2: Observing the Context}
	
	\textbf{Assessment Criteria:} C2
	
	\textbf{Learning Purpose:}
	
	Understand a potential problem or opportunity through direct and systematic observation.
	
	\textbf{Main Activity:}
	
	Teams select one opportunity from their initial exploration and conduct a structured observation of the relevant context.
	
	Students observe spaces, behaviors, users, interactions, and situations related to the selected opportunity. They record what they see and identify recurring patterns and contextual factors.
	
	Students are encouraged to distinguish between what they directly observe and what they assume or interpret.
	
	\textbf{Student Product:}
	
	Observation record containing relevant observations, contextual information, and emerging patterns.
	
	\textbf{Formative Assessment:}
	
	The teacher reviews the observation records and provides feedback on the quality and relevance of the information collected.
	
	% -------------------------------------------------
	% 7.3 CLASS 3
	% -------------------------------------------------
	
	\subsection{Class 3: Understanding the User}
	
	\textbf{Assessment Criteria:} C3--C4
	
	\textbf{Learning Purpose:}
	
	Understand the experiences and perspectives of people affected by the identified situation and distinguish real needs from anticipated solutions.
	
	\textbf{Main Activity:}
	
	Students identify potential users or people affected by the situation being investigated.
	
	Teams prepare a small set of guided questions and conduct short interviews, conversations, or surveys. Students collect information about users' experiences, needs, frustrations, expectations, and perceptions.
	
	The information collected is compared with the observations from Class 2.
	
	Students analyze whether the evidence supports the original problem definition or suggests that it should be reformulated.
	
	\textbf{Student Product:}
	
	User profile and summary of user findings.
	
	\textbf{Formative Assessment:}
	
	The teacher reviews whether students are listening to users and whether their conclusions are supported by information obtained from them.
	
	% -------------------------------------------------
	% 7.4 CLASS 4
	% -------------------------------------------------
	
	\subsection{Class 4: Building the Contextual Diagnosis}
	
	\textbf{Assessment Criteria:} C1--C5
	
	\textbf{Learning Purpose:}
	
	Organize the information collected during the investigation and construct an evidence-based contextual diagnosis.
	
	\textbf{Main Activity:}
	
	Teams consolidate the information collected through observation, interviews, conversations, surveys, and contextual exploration.
	
	Students organize their findings using appropriate representations such as:
	
	\begin{itemize}
		\item Context maps.
		\item User profiles.
		\item Empathy maps.
		\item Lists of findings.
		\item Cause-and-consequence diagrams.
		\item Simple system or process diagrams.
	\end{itemize}
	
	Students identify the main findings and analyze the causes, consequences, users affected, and relevance of the situation.
	
	They must explicitly distinguish the identified need or problem from a possible solution.
	
	\textbf{Student Product:}
	
	\textbf{Learning Evidence 1: Contextual Diagnosis.}
	
	The evidence should demonstrate the development of Criteria C1--C5.
	
	\textbf{Assessment:}
	
	Learning Evidence 1 is assessed according to the students' ability to identify a relevant opportunity, understand the context, listen to users, distinguish needs from solutions, and organize contextual information.
	
	% -------------------------------------------------
	% LEARNING EVIDENCE 1
	% -------------------------------------------------
	
	\subsection{Learning Evidence 1: Contextual Diagnosis}
	
	\textbf{Assessment Criteria:} C1--C5
	
	Students work in teams to investigate a problem, need, discomfort, or opportunity within the school or nearby community.
	
	The evidence should contain:
	
	\begin{enumerate}
		\item An identified need, problem, or opportunity.
		\item A description of the relevant context.
		\item Observation records.
		\item Identification of relevant users.
		\item A user profile or empathy map.
		\item Interview, conversation, or survey findings.
		\item Main findings and relevant patterns.
		\item Causes and consequences of the situation.
		\item Evidence supporting the existence and relevance of the problem or opportunity.
		\item A context map, diagram, or other visual organization of the information.
		\item An explicit distinction between the identified need and a possible solution.
	\end{enumerate}
	
	\textbf{Central Question:}
	
	\begin{quote}
		What is actually happening, who is affected, and what evidence do we have that this is a relevant need or opportunity?
	\end{quote}
	
	The evidence should remain primarily diagnostic rather than solution-oriented.
	
	% -------------------------------------------------
	% 7.5 CLASS 5
	% -------------------------------------------------
	
	\subsection{Class 5: Prioritizing the Innovation Opportunity}
	
	\textbf{Assessment Criteria:} C6
	
	\textbf{Learning Purpose:}
	
	Select an innovation opportunity by considering its relevance, feasibility, and potential benefit.
	
	\textbf{Main Activity:}
	
	Teams review the findings from Learning Evidence 1 and identify the opportunities that could potentially be addressed.
	
	Students compare the opportunities using criteria such as:
	
	\begin{itemize}
		\item Relevance to the users or community.
		\item Magnitude or importance of the problem.
		\item Feasibility of taking action.
		\item Availability of resources.
		\item Potential benefit or positive impact.
	\end{itemize}
	
	Teams organize their comparison using a simple prioritization matrix or scoring process.
	
	They then select one opportunity and justify their decision using evidence from the contextual diagnosis.
	
	\textbf{Student Product:}
	
	Opportunity prioritization matrix and written justification.
	
	\textbf{Formative Assessment:}
	
	The teacher verifies that the selected opportunity is supported by the evidence collected during the previous classes and that the decision is based on explicit criteria rather than personal preference.
	
	% -------------------------------------------------
	% 7.6 CLASS 6
	% -------------------------------------------------
	
	\subsection{Class 6: Defining the Innovation Challenge}
	
	\textbf{Assessment Criteria:} C7--C8, C10
	
	\textbf{Learning Purpose:}
	
	Transform the selected opportunity into a clear, concrete, and understandable innovation challenge while maintaining a strong focus on the user.
	
	\textbf{Main Activity:}
	
	Teams review the selected opportunity and the evidence gathered in Learning Evidence 1.
	
	They identify:
	
	\begin{itemize}
		\item The main user or beneficiary.
		\item The specific situation that needs attention.
		\item The relevant need or problem.
		\item The desired direction of improvement.
	\end{itemize}
	
	Students formulate an innovation challenge that is sufficiently specific to guide future ideation while remaining open enough to allow multiple possible solutions.
	
	The challenge is then connected with knowledge from one or more relevant disciplines, such as:
	
	\begin{itemize}
		\item Design.
		\item Communication.
		\item Economics.
		\item Technology.
		\item Citizenship.
		\item Environmental studies.
		\item Social sciences.
		\item Other relevant areas.
	\end{itemize}
	
	Teams exchange their challenge statements and provide peer feedback regarding clarity, relevance, evidence, and user focus.
	
	\textbf{Student Product:}
	
	Innovation challenge statement and disciplinary connection.
	
	\textbf{Formative Assessment:}
	
	The teacher reviews whether the challenge is clear, concrete, evidence-based, connected to a relevant user, and sufficiently open to multiple possible solutions.
	
	% -------------------------------------------------
	% 7.7 CLASS 7
	% -------------------------------------------------
	
	\subsection{Class 7: Initial Ideation}
	
	\textbf{Assessment Criteria:} C9--C10
	
	\textbf{Learning Purpose:}
	
	Explore possible directions for addressing the innovation challenge while remaining open to diverse proposals and maintaining a user-centered perspective.
	
	\textbf{Main Activity:}
	
	Students begin with the innovation challenge defined in Class 6.
	
	Teams participate in divergent ideation activities designed to generate multiple possible approaches before selecting a specific solution.
	
	Possible activities include:
	
	\begin{itemize}
		\item Brainstorming.
		\item ``How might we...?'' questions.
		\item Individual idea generation.
		\item Team idea generation.
		\item Rapid sketching.
		\item Idea clustering.
		\item Comparison of alternative approaches.
	\end{itemize}
	
	Students are encouraged to generate a variety of possibilities rather than immediately selecting the first feasible idea.
	
	At the end of the class, teams organize their initial ideas and explain how these ideas relate to the identified user and innovation challenge.
	
	\textbf{Student Product:}
	
	\textbf{Learning Evidence 2: Innovation Challenge and Initial Ideation.}
	
	\textbf{Assessment:}
	
	Learning Evidence 2 is assessed according to Criteria C6--C10.
	
	% -------------------------------------------------
	% LEARNING EVIDENCE 2
	% -------------------------------------------------
	
	\subsection{Learning Evidence 2: Innovation Challenge and Initial Ideation}
	
	\textbf{Assessment Criteria:} C6--C10
	
	Using Learning Evidence 1 as its starting point, students select the most relevant opportunity and transform it into a clearly defined innovation challenge. They then explore possible approaches through initial ideation.
	
	The evidence should contain:
	
	\begin{enumerate}
		\item The selected innovation opportunity.
		\item The prioritization criteria and justification.
		\item The relevance of the selected opportunity.
		\item Feasibility considerations.
		\item Expected benefit or potential impact.
		\item A clearly formulated innovation challenge.
		\item Identification of the principal user or beneficiary.
		\item Connection with at least one relevant disciplinary area.
		\item An initial set of possible ideas or approaches.
		\item Evidence of divergent thinking and openness to different proposals.
		\item An explanation of why understanding the user is necessary before developing a final solution.
	\end{enumerate}
	
	\textbf{Central Question:}
	
	\begin{quote}
		Given what we have learned about the context and users, which opportunity should we address, what challenge will we pursue, and what possible directions could lead to an innovative solution?
	\end{quote}
	
	% -------------------------------------------------
	% 7.8 PEDAGOGICAL PROGRESSION
	% -------------------------------------------------
	
	\subsection{Pedagogical Progression}
	
	The seven classes follow a progressive sequence from contextual exploration to initial ideation.
	
	\begin{center}
		\begin{tabular}{c}
			\textbf{Class 1: Identify} [0.1cm]
			$\downarrow$ [0.1cm]
			\textbf{Class 2: Observe} [0.1cm]
			$\downarrow$ [0.1cm]
			\textbf{Class 3: Listen to Users} [0.1cm]
			$\downarrow$ [0.1cm]
			\textbf{Class 4: Organize and Diagnose} [0.1cm]
			$\downarrow$ [0.15cm]
			\boxed{\textbf{Learning Evidence 1: Contextual Diagnosis}} [-0.05cm]
			\boxed{\textbf{C1--C5}} [0.2cm]
			$\downarrow$ [0.1cm]
			\textbf{Class 5: Prioritize} [0.1cm]
			$\downarrow$ [0.1cm]
			\textbf{Class 6: Define and Connect} [0.1cm]
			$\downarrow$ [0.1cm]
			\textbf{Class 7: Ideate} [0.1cm]
			$\downarrow$ [0.15cm]
			\boxed{\textbf{Learning Evidence 2: Innovation Challenge and Initial Ideation}} [-0.05cm]
			\boxed{\textbf{C6--C10}}
		\end{tabular}
	\end{center}
	
	% -------------------------------------------------
	% 7.9 ALIGNMENT
	% -------------------------------------------------
	
	\subsection{Alignment Between Classes, Criteria, and Evidence}
	
	\subsubsection*{Learning Evidence 1}
	
	\textbf{Classes:} 1--4
	
	\textbf{Criteria:} C1--C5
	
	\textbf{Focus:}
	
	Contextual exploration, observation, user understanding, distinction between needs and solutions, and organization of contextual information.
	
	\textbf{Final Product:}
	
	Contextual Diagnosis.
	
	\subsubsection*{Learning Evidence 2}
	
	\textbf{Classes:} 5--7
	
	\textbf{Criteria:} C6--C10
	
	\textbf{Focus:}
	
	Opportunity prioritization, innovation challenge definition, disciplinary connection, user-centered thinking, and initial ideation.
	
	\textbf{Final Product:}
	
	Innovation Challenge and Initial Ideation.
	
	% -------------------------------------------------
	% 7.10 SCOPE
	% -------------------------------------------------
	
	\subsection{Scope of the Seven-Class Experience}
	
	The seven-class sequence is intentionally focused on identifying and understanding problems and innovation opportunities rather than developing a complete technological or entrepreneurial solution.
	
	Therefore, the expected outcome at the end of the sequence is not a finished product, application, business model, or prototype.
	
	Instead, students should have:
	
	\begin{itemize}
		\item Identified and investigated a relevant need or opportunity.
		\item Collected evidence from the context and potential users.
		\item Organized the evidence into a contextual diagnosis.
		\item Distinguished a real need from an anticipated solution.
		\item Prioritized one opportunity using explicit criteria.
		\item Defined a clear innovation challenge.
		\item Connected the challenge with relevant disciplinary knowledge.
		\item Generated initial ideas and possible directions for a future solution.
	\end{itemize}
	
	Subsequent learning experiences may address prototyping, validation, budgeting, business models, implementation, impact evaluation, and refinement of the proposed solution.
	
	
	% =================================================
	% 7.1 PAUSE
	% =================================================
	
	\subsection{7.1 Pause: What Am I Achieving?}
	
	\textbf{Description:}
	
	Description focused on ``What am I achieving?''
	
	\vspace{0.3cm}
	
	\textbf{Way of Proceeding:}
	
	\placeholder{To be defined.}
	
	\textbf{Way of Being:}
	
	\placeholder{To be defined.}
	
	% =================================================
	% 7.2 PAUSE
	% =================================================
	
	\subsection{7.2 Pause: What Did I Achieve?}
	
	\textbf{Description:}
	
	Description focused on ``What did I achieve?''
	
	\vspace{0.3cm}
	
	\textbf{Way of Proceeding:}
	
	\placeholder{To be defined.}
	
	\textbf{Way of Being:}
	
	\placeholder{To be defined.}
	
	% =================================================
	% REFERENCES
	% =================================================
	
	\section*{References}
	
	\placeholder{To be completed.}
	
\end{document}